\Chapter{73}

\Verse{1}
{
    \zA{话说那赵姨娘和贾政说话，忽听外面一声响，不知何物，忙问时，原来是外间 窗屉不曾扣好，滑了屈戌掉下来。}
}
{
    \eA{[English source not present in the checked-in Bencraft/Joly files.]}
}
{
}

\Verse{2}
{
    \zA{赵姨娘骂了丫头几句，自己带领丫鬟上好，方进 来打发贾政安歇，不在话下。}
}
{
    \eA{[English source not present in the checked-in Bencraft/Joly files.]}
}
{
}

\Verse{3}
{
    \zA{却说怡红院中宝玉方才睡下，丫鬟们正欲各散安歇，忽听有人来敲院门。}
}
{
    \eA{[English source not present in the checked-in Bencraft/Joly files.]}
}
{
}

\Verse{4}
{
    \zA{老婆 子开了，见是赵姨娘房内的丫头名唤小鹊的，问他作什么，小鹊不答，直往里走， 来找宝玉。}
}
{
    \eA{[English source not present in the checked-in Bencraft/Joly files.]}
}
{
}

\Verse{5}
{
    \zA{只见宝玉才睡下，晴雯等犹在床边坐着，大家玩笑。}
}
{
    \eA{[English source not present in the checked-in Bencraft/Joly files.]}
}
{
}

\Verse{6}
{
    \zA{见他来了，都问： “什么事，这时候又跑了来？”}
}
{
    \eA{[English source not present in the checked-in Bencraft/Joly files.]}
}
{
}

\Verse{7}
{
    \zA{小鹊连忙悄向宝玉道：“我来告诉你个信儿，方才 我们奶奶咕咕唧唧的，在老爷前不知说了你些个什么，我只听见‘宝玉’二字。}
}
{
    \eA{[English source not present in the checked-in Bencraft/Joly files.]}
}
{
}

\Verse{8}
{
    \zA{我 来告诉你，仔细明儿老爷和你说话罢。”}
}
{
    \eA{[English source not present in the checked-in Bencraft/Joly files.]}
}
{
}

\Verse{9}
{
    \zA{一面说着，回身就走。}
}
{
    \eA{[English source not present in the checked-in Bencraft/Joly files.]}
}
{
}

\Verse{10}
{
    \zA{袭人命人留他吃茶， 因怕关门，遂一直去了。}
}
{
    \eA{[English source not present in the checked-in Bencraft/Joly files.]}
}
{
}

\Verse{11}
{
    \zA{宝玉听了，知道赵姨娘心术不端，合自己仇人似的，又不 知他说些什么，便如孙大圣听见了紧箍儿咒的一般，登时四肢五内一齐皆不自在起 来。}
}
{
    \eA{[English source not present in the checked-in Bencraft/Joly files.]}
}
{
}

\Verse{12}
{
    \zA{想来想去，别无他法，且理熟了书预备明儿盘考，只能书不舛错，就有别事也 可搪塞。}
}
{
    \eA{[English source not present in the checked-in Bencraft/Joly files.]}
}
{
}

\Verse{13}
{
    \zA{一面想罢，忙披衣起来要读书。}
}
{
    \eA{[English source not present in the checked-in Bencraft/Joly files.]}
}
{
}

\Verse{14}
{
    \zA{心中又自后悔：“这些日子，只说不提了， 偏又丢生了。}
}
{
    \eA{[English source not present in the checked-in Bencraft/Joly files.]}
}
{
}

\Verse{15}
{
    \zA{早知该天天好歹温习些。”}
}
{
    \eA{[English source not present in the checked-in Bencraft/Joly files.]}
}
{
}

\Verse{16}
{
    \zA{如今打算打算，肚子里现可背诵的，不过 只有《学》、《庸》、二《论》还背得出来。}
}
{
    \eA{[English source not present in the checked-in Bencraft/Joly files.]}
}
{
}

\Verse{17}
{
    \zA{至上本《孟子》，就有一半是夹生的， 若凭空提一句，断不能背；至下《孟子》，就有大半生的，算起《五经》来，因近
        来做诗，常把《五经》集些，虽不甚熟，还可塞责。}
}
{
    \eA{[English source not present in the checked-in Bencraft/Joly files.]}
}
{
}

\Verse{18}
{
    \zA{别的虽不记得，素日贾政幸未 叫读的，纵不知，也还不妨。}
}
{
    \eA{[English source not present in the checked-in Bencraft/Joly files.]}
}
{
}

\Verse{19}
{
    \zA{至于古文，还是那几年所读过的几篇《左传》、《国 策》、《公羊》、《谷梁》、汉、唐等文，这几年未曾读得，不过一时之兴，随看 随忘，未曾下过苦功，如何记得？}
}
{
    \eA{[English source not present in the checked-in Bencraft/Joly files.]}
}
{
}

\Verse{20}
{
    \zA{这是更难塞责的。}
}
{
    \eA{[English source not present in the checked-in Bencraft/Joly files.]}
}
{
}

\Verse{21}
{
    \zA{更有时文八股一道，因平素深 恶，说这原非圣贤之制撰，焉能阐发圣贤之奥，不过是后人饵名钓禄之阶。}
}
{
    \eA{[English source not present in the checked-in Bencraft/Joly files.]}
}
{
}

\Verse{22}
{
    \zA{虽贾政 当日起身，选了百十篇命他读的，不过是后人的时文，偶见其中一二股内，或承起
        之中，有作的精致，或流荡、或游戏、或悲感稍能动性者，偶尔一读，不过供一时 之兴趣，究竟何曾成篇潜心玩索？}
}
{
    \eA{[English source not present in the checked-in Bencraft/Joly files.]}
}
{
}

\Verse{23}
{
    \zA{如今若温习这个，又恐明日盘究那个；若温习那 个，又恐盘驳这个：一夜之工，亦不能全然温习。}
}
{
    \eA{[English source not present in the checked-in Bencraft/Joly files.]}
}
{
}

\Verse{24}
{
    \zA{因此，越添了焦躁。}
}
{
    \eA{[English source not present in the checked-in Bencraft/Joly files.]}
}
{
}

\Verse{25}
{
    \zA{自己读书，不值紧要，却累着一房丫鬟们都不能睡。}
}
{
    \eA{[English source not present in the checked-in Bencraft/Joly files.]}
}
{
}

\Verse{26}
{
    \zA{袭人等在旁剪烛斟茶，那 些小的都困倦起来，前仰后合。}
}
{
    \eA{[English source not present in the checked-in Bencraft/Joly files.]}
}
{
}

\Verse{27}
{
    \zA{晴雯骂道：“什么小蹄子们!一个个黑家白日挺尸 挺不够，偶然一次睡迟了些，就装出这个腔调儿来了。}
}
{
    \eA{[English source not present in the checked-in Bencraft/Joly files.]}
}
{
}

\Verse{28}
{
    \zA{再这么着，我拿针扎你们两 下子！”}
}
{
    \eA{[English source not present in the checked-in Bencraft/Joly files.]}
}
{
}

\Verse{29}
{
    \zA{话犹未了，只听外间咕咚一声，急忙看时，原来是个小丫头坐着打盹，一 头撞到壁上，从梦中惊醒。}
}
{
    \eA{[English source not present in the checked-in Bencraft/Joly files.]}
}
{
}

\Verse{30}
{
    \zA{却正是晴雯说这话之时，他怔怔的只当是晴雯打了他一 下子，遂哭着央说：“好姐姐，我再不敢了！”}
}
{
    \eA{[English source not present in the checked-in Bencraft/Joly files.]}
}
{
}

\Verse{31}
{
    \zA{众人都笑起来。}
}
{
    \eA{[English source not present in the checked-in Bencraft/Joly files.]}
}
{
}

\Verse{32}
{
    \zA{宝玉忙劝道：“饶 他罢。}
}
{
    \eA{[English source not present in the checked-in Bencraft/Joly files.]}
}
{
}

\Verse{33}
{
    \zA{原该叫他们睡去。}
}
{
    \eA{[English source not present in the checked-in Bencraft/Joly files.]}
}
{
}

\Verse{34}
{
    \zA{你们也该替换着睡。”}
}
{
    \eA{[English source not present in the checked-in Bencraft/Joly files.]}
}
{
}

\Verse{35}
{
    \zA{袭人道：“小祖宗，你只顾你的罢! 统共这一夜的工夫，你把心暂且用在这几本书上，等过了这一关，由你再张罗别的， 也不算误了什么。”}
}
{
    \eA{[English source not present in the checked-in Bencraft/Joly files.]}
}
{
}

\Verse{36}
{
    \zA{宝玉听他说的恳切，只得又读几句。}
}
{
    \eA{[English source not present in the checked-in Bencraft/Joly files.]}
}
{
}

\Verse{37}
{
    \zA{麝月斟了一杯茶来润舌， 宝玉接茶吃了。}
}
{
    \eA{[English source not present in the checked-in Bencraft/Joly files.]}
}
{
}

\Verse{38}
{
    \zA{因见麝月只穿着短袄，宝玉道：“夜静了冷，到底穿一件大衣裳才 是啊。”}
}
{
    \eA{[English source not present in the checked-in Bencraft/Joly files.]}
}
{
}

\Verse{39}
{
    \zA{麝月笑指着书道：“你暂且把我们忘了，使不得吗？}
}
{
    \eA{[English source not present in the checked-in Bencraft/Joly files.]}
}
{
}

\Verse{40}
{
    \zA{且把心搁在这上头些 罢。”}
}
{
    \eA{[English source not present in the checked-in Bencraft/Joly files.]}
}
{
}

\Verse{41}
{
    \zA{话犹未了，只听春燕秋纹从后房门跑进来，口内喊说：“不好了!一个人打墙 上跳下来了。”}
}
{
    \eA{[English source not present in the checked-in Bencraft/Joly files.]}
}
{
}

\Verse{42}
{
    \zA{众人听说，忙问：“在那里？”}
}
{
    \eA{[English source not present in the checked-in Bencraft/Joly files.]}
}
{
}

\Verse{43}
{
    \zA{即喝起人来，各处寻找。}
}
{
    \eA{[English source not present in the checked-in Bencraft/Joly files.]}
}
{
}

\Verse{44}
{
    \zA{晴雯因见 宝玉读书苦恼，劳费一夜神思，明日也未必妥当，心下正要替宝玉想个主意，好脱 此难。}
}
{
    \eA{[English source not present in the checked-in Bencraft/Joly files.]}
}
{
}

\Verse{45}
{
    \zA{忽然碰着这一惊，便生计向宝玉道：“趁这个机会，快装病，只说吓着了。”}
}
{
    \eA{[English source not present in the checked-in Bencraft/Joly files.]}
}
{
}

\Verse{46}
{
    \zA{这话正中宝玉心怀。}
}
{
    \eA{[English source not present in the checked-in Bencraft/Joly files.]}
}
{
}

\Verse{47}
{
    \zA{因叫起上夜的来，打着灯笼各处搜寻，并无踪迹，都说：“小 姑娘们想是睡花了眼出去，风摇的树枝儿，错认了人。”}
}
{
    \eA{[English source not present in the checked-in Bencraft/Joly files.]}
}
{
}

\Verse{48}
{
    \zA{晴雯便道：“别放屁!你 们查的不严，怕耽不是，还拿这话来支吾!刚才并不是一个人见的，宝玉和我们出 去，大家亲见的。}
}
{
    \eA{[English source not present in the checked-in Bencraft/Joly files.]}
}
{
}

\Verse{49}
{
    \zA{如今宝玉吓得颜色都变了，满身发热，我这会子还要上房里取安 魂丸药去呢。}
}
{
    \eA{[English source not present in the checked-in Bencraft/Joly files.]}
}
{
}

\Verse{50}
{
    \zA{太太问起来，是要回明白了的，难道依你说就罢了？”}
}
{
    \eA{[English source not present in the checked-in Bencraft/Joly files.]}
}
{
}

\Verse{51}
{
    \zA{众人听了吓得 不敢则声，只得又各处去找。}
}
{
    \eA{[English source not present in the checked-in Bencraft/Joly files.]}
}
{
}

\Verse{52}
{
    \zA{晴雯和秋纹二人果出去要药去，故意闹的众人皆知宝 玉着了惊，吓病了。}
}
{
    \eA{[English source not present in the checked-in Bencraft/Joly files.]}
}
{
}

\Verse{53}
{
    \zA{王夫人听了，忙命人来看视给药，又吩咐各上夜人仔细搜查； 又一面叫查二门外邻园墙上夜的小厮们。}
}
{
    \eA{[English source not present in the checked-in Bencraft/Joly files.]}
}
{
}

\Verse{54}
{
    \zA{于是园内灯笼火把，直闹了一夜。}
}
{
    \eA{[English source not present in the checked-in Bencraft/Joly files.]}
}
{
}

\Verse{55}
{
    \zA{至五更 天，就传管家的细看查访。}
}
{
    \eA{[English source not present in the checked-in Bencraft/Joly files.]}
}
{
}

\Verse{56}
{
    \zA{贾母闻知宝玉被吓，细问原由，众人不敢再隐，只得回明。}
}
{
    \eA{[English source not present in the checked-in Bencraft/Joly files.]}
}
{
}

\Verse{57}
{
    \zA{贾母道：“我不料 有此事。}
}
{
    \eA{[English source not present in the checked-in Bencraft/Joly files.]}
}
{
}

\Verse{58}
{
    \zA{如今各处上夜的都不小心还是小事，只怕他们就是贼也未可知。”}
}
{
    \eA{[English source not present in the checked-in Bencraft/Joly files.]}
}
{
}

\Verse{59}
{
    \zA{当下邢 夫人尤氏等都过来请安，李纨凤姐及姊妹等皆陪侍，听贾母如此说，都默无所答。}
}
{
    \eA{[English source not present in the checked-in Bencraft/Joly files.]}
}
{
}

\Verse{60}
{
    \zA{独探春出位笑道：“近因凤姐姐身子不好几日，园里的人，比先放肆许多。}
}
{
    \eA{[English source not present in the checked-in Bencraft/Joly files.]}
}
{
}

\Verse{61}
{
    \zA{先前不 过是大家偷着一时半刻，或夜里坐更时三四个人聚在一处，或掷骰，或斗牌，小玩 意儿，不过为着熬困起见。}
}
{
    \eA{[English source not present in the checked-in Bencraft/Joly files.]}
}
{
}

\Verse{62}
{
    \zA{如今渐次放诞，竟开了赌局，甚至头家局主，或三十吊 五十吊的大输赢。}
}
{
    \eA{[English source not present in the checked-in Bencraft/Joly files.]}
}
{
}

\Verse{63}
{
    \zA{半月前竟有争斗相打的事。”}
}
{
    \eA{[English source not present in the checked-in Bencraft/Joly files.]}
}
{
}

\Verse{64}
{
    \zA{贾母听了，忙说：“你既知道，为 什么不早回我来？”}
}
{
    \eA{[English source not present in the checked-in Bencraft/Joly files.]}
}
{
}

\Verse{65}
{
    \zA{探春道：“我因想着太太事多，且连日不自在，所以没回，只 告诉大嫂子和管事的人们，戒饬过几次，近日好些了。”}
}
{
    \eA{[English source not present in the checked-in Bencraft/Joly files.]}
}
{
}

\Verse{66}
{
    \zA{贾母忙道：“你姑娘家， 那里知道这里头的利害？}
}
{
    \eA{[English source not present in the checked-in Bencraft/Joly files.]}
}
{
}

\Verse{67}
{
    \zA{你以为赌钱常事，不过怕起争端；不知夜间既耍钱，就保 不住不吃酒，既吃酒，就未免门户任意开锁，或买东西，其中夜静人稀，趁便藏贼 引盗，什么事做不出来？}
}
{
    \eA{[English source not present in the checked-in Bencraft/Joly files.]}
}
{
}

\Verse{68}
{
    \zA{况且园内你姐儿们起居所伴者，皆系丫头媳妇们，贤愚混 杂。}
}
{
    \eA{[English source not present in the checked-in Bencraft/Joly files.]}
}
{
}

\Verse{69}
{
    \zA{贼盗事小，倘有别事，略沾带些，关系非小!这事岂可轻恕？”}
}
{
    \eA{[English source not present in the checked-in Bencraft/Joly files.]}
}
{
}

\Verse{70}
{
    \zA{探春听说，便 默然归坐。}
}
{
    \eA{[English source not present in the checked-in Bencraft/Joly files.]}
}
{
}

\Verse{71}
{
    \zA{凤姐虽未大愈，精神未尝稍减，今见贾母如此说，便忙道：“偏偏我又 病了。”}
}
{
    \eA{[English source not present in the checked-in Bencraft/Joly files.]}
}
{
}

\Verse{72}
{
    \zA{遂回头命人速传林之孝家的等总理家事的四个媳妇来了，当着贾母申饬了 一顿。}
}
{
    \eA{[English source not present in the checked-in Bencraft/Joly files.]}
}
{
}

\Verse{73}
{
    \zA{贾母命：“即刻查了头家赌家来!有人出首者赏，隐情不告者罚。”}
}
{
    \eA{[English source not present in the checked-in Bencraft/Joly files.]}
}
{
}

\Verse{74}
{
    \zA{林之孝家的等见贾母动怒，谁敢徇私，忙去园内传齐，又一一盘查。}
}
{
    \eA{[English source not present in the checked-in Bencraft/Joly files.]}
}
{
}

\Verse{75}
{
    \zA{虽然大家 赖一回，终不免水落石出。}
}
{
    \eA{[English source not present in the checked-in Bencraft/Joly files.]}
}
{
}

\Verse{76}
{
    \zA{查得大头家三人，小头家八人，聚赌者统共二十多人， 都带来见贾母，跪在院内，磕响头求饶。}
}
{
    \eA{[English source not present in the checked-in Bencraft/Joly files.]}
}
{
}

\Verse{77}
{
    \zA{贾母先问大头家名姓，和钱之多少。}
}
{
    \eA{[English source not present in the checked-in Bencraft/Joly files.]}
}
{
}

\Verse{78}
{
    \zA{原来 这大头家，一个是林之孝家的两姨亲家，一个是园内厨房内柳家媳妇之妹，一个是 迎春之乳母。}
}
{
    \eA{[English source not present in the checked-in Bencraft/Joly files.]}
}
{
}

\Verse{79}
{
    \zA{这是三个为首的，馀者不能多记。}
}
{
    \eA{[English source not present in the checked-in Bencraft/Joly files.]}
}
{
}

\Verse{80}
{
    \zA{贾母便命将骰子纸牌一并烧毁，所 有的钱入官，分散与众人；将为首者每人打四十大板，撵出去，总不许再入；从者 每人打二十板，革去三月月钱，拨入圊厕行内。}
}
{
    \eA{[English source not present in the checked-in Bencraft/Joly files.]}
}
{
}

\Verse{81}
{
    \zA{又将林之孝家的申饬了一番。}
}
{
    \eA{[English source not present in the checked-in Bencraft/Joly files.]}
}
{
}

\Verse{82}
{
    \zA{林之 孝家的见他的亲戚又给他打嘴，自己也觉没趣；迎春在坐也觉没意思。}
}
{
    \eA{[English source not present in the checked-in Bencraft/Joly files.]}
}
{
}

\Verse{83}
{
    \zA{黛玉、宝钗、 探春等见迎春的乳母如此，也是“物伤其类”的意思，遂都起身笑向贾母讨情，说：
        “这个奶奶素日原不玩的，不知怎么，也偶然高兴；求看二姐姐面上，饶过这次罢。”}
}
{
    \eA{[English source not present in the checked-in Bencraft/Joly files.]}
}
{
}

\Verse{84}
{
    \zA{贾母道：“你们不知道。}
}
{
    \eA{[English source not present in the checked-in Bencraft/Joly files.]}
}
{
}

\Verse{85}
{
    \zA{大约这些奶子们，一个个仗着奶过哥儿姐儿，原比别人有 些体面，他们就生事，比别人更可恶!专管调唆主子，护短偏向。}
}
{
    \eA{[English source not present in the checked-in Bencraft/Joly files.]}
}
{
}
\ChapterImage{img/chapters/141第73-74回 癡丫頭誤拾繡春囊 懦小姐不問累金鳳.jpg}


\Verse{86}
{
    \zA{我都是经过的。}
}
{
    \eA{[English source not present in the checked-in Bencraft/Joly files.]}
}
{
}

\Verse{87}
{
    \zA{况且要拿一个作法，恰好果然就遇见了一个。}
}
{
    \eA{[English source not present in the checked-in Bencraft/Joly files.]}
}
{
}

\Verse{88}
{
    \zA{你们别管，我自有道理。”}
}
{
    \eA{[English source not present in the checked-in Bencraft/Joly files.]}
}
{
}

\Verse{89}
{
    \zA{宝钗等听 说，只得罢了。}
}
{
    \eA{[English source not present in the checked-in Bencraft/Joly files.]}
}
{
}

\Verse{90}
{
    \zA{一时贾母歇晌，大家散出，都知贾母生气，皆不敢回家，只得在此 暂候。}
}
{
    \eA{[English source not present in the checked-in Bencraft/Joly files.]}
}
{
}

\Verse{91}
{
    \zA{尤氏到凤姐儿处来闲话了一回，因他也不自在，只得园内去闲谈。}
}
{
    \eA{[English source not present in the checked-in Bencraft/Joly files.]}
}
{
}

\Verse{92}
{
    \zA{邢夫人在王夫人处坐了一回，也要到园内走走。}
}
{
    \eA{[English source not present in the checked-in Bencraft/Joly files.]}
}
{
}

\Verse{93}
{
    \zA{刚至园门前，只见贾母房内的 小丫头子名唤傻大姐的，笑嘻嘻走来，手内拿着个花红柳绿的东西，低头瞧着只管 走。}
}
{
    \eA{[English source not present in the checked-in Bencraft/Joly files.]}
}
{
}

\Verse{94}
{
    \zA{不防迎头撞见邢夫人，抬头看见，方才站住。}
}
{
    \eA{[English source not present in the checked-in Bencraft/Joly files.]}
}
{
}

\Verse{95}
{
    \zA{邢夫人因说：“这傻丫头又得个 什么爱巴物儿，这样喜欢？}
}
{
    \eA{[English source not present in the checked-in Bencraft/Joly files.]}
}
{
}

\Verse{96}
{
    \zA{拿来我瞧瞧。”}
}
{
    \eA{[English source not present in the checked-in Bencraft/Joly files.]}
}
{
}

\Verse{97}
{
    \zA{原来这傻大姐年方十四岁，是新挑上来 给贾母这边专做粗活的。}
}
{
    \eA{[English source not present in the checked-in Bencraft/Joly files.]}
}
{
}

\Verse{98}
{
    \zA{因他生的体肥面阔，两只大脚，做粗活很爽利简捷，且心 性愚顽，一无知识，出言可以发笑。}
}
{
    \eA{[English source not present in the checked-in Bencraft/Joly files.]}
}
{
}

\Verse{99}
{
    \zA{贾母喜欢，便起名为“傻大姐”，若有错失， 也不苛责他。}
}
{
    \eA{[English source not present in the checked-in Bencraft/Joly files.]}
}
{
}

\Verse{100}
{
    \zA{无事时便入园内来玩耍，正往山石背后掏促织去，忽见一个五彩绣香 囊，上面绣的并非花鸟等物，一面却是两个人赤条条的相抱，一面是几个字。}
}
{
    \eA{[English source not present in the checked-in Bencraft/Joly files.]}
}
{
}

\Verse{101}
{
    \zA{这痴 丫头原不认得是春意儿，心下打量：“敢是两个妖精打架？}
}
{
    \eA{[English source not present in the checked-in Bencraft/Joly files.]}
}
{
}

\Verse{102}
{
    \zA{不就是两个人打架呢？”}
}
{
    \eA{[English source not present in the checked-in Bencraft/Joly files.]}
}
{
}

\Verse{103}
{
    \zA{左右猜解不来，正要拿去给贾母看呢，所以笑嘻嘻走回。}
}
{
    \eA{[English source not present in the checked-in Bencraft/Joly files.]}
}
{
}

\Verse{104}
{
    \zA{忽见邢夫人如此说，便笑 道：“太太真个说的巧，真是个爱巴物儿。}
}
{
    \eA{[English source not present in the checked-in Bencraft/Joly files.]}
}
{
}

\Verse{105}
{
    \zA{太太瞧一瞧。”}
}
{
    \eA{[English source not present in the checked-in Bencraft/Joly files.]}
}
{
}

\Verse{106}
{
    \zA{说着便送过去。}
}
{
    \eA{[English source not present in the checked-in Bencraft/Joly files.]}
}
{
}

\Verse{107}
{
    \zA{邢夫人 接来一看，吓得连忙死紧攥住，忙问：“你是那里得的？”}
}
{
    \eA{[English source not present in the checked-in Bencraft/Joly files.]}
}
{
}

\Verse{108}
{
    \zA{傻大姐道：“我掏促织 儿，在山子石后头拣的。”}
}
{
    \eA{[English source not present in the checked-in Bencraft/Joly files.]}
}
{
}

\Verse{109}
{
    \zA{邢夫人道：“快别告诉人!这不是好东西。}
}
{
    \eA{[English source not present in the checked-in Bencraft/Joly files.]}
}
{
}

\Verse{110}
{
    \zA{连你也要打 死呢。}
}
{
    \eA{[English source not present in the checked-in Bencraft/Joly files.]}
}
{
}

\Verse{111}
{
    \zA{因你素日是个傻丫头，以后再别提了。”}
}
{
    \eA{[English source not present in the checked-in Bencraft/Joly files.]}
}
{
}

\Verse{112}
{
    \zA{这傻大姐听了，反吓得黄了脸，说： “再不敢了。”}
}
{
    \eA{[English source not present in the checked-in Bencraft/Joly files.]}
}
{
}

\Verse{113}
{
    \zA{磕了头，呆呆而去。}
}
{
    \eA{[English source not present in the checked-in Bencraft/Joly files.]}
}
{
}

\Verse{114}
{
    \zA{邢夫人回头看时，都是些女孩儿，不便递给他们，自己便？}
}
{
    \eA{[English source not present in the checked-in Bencraft/Joly files.]}
}
{
}

\Verse{115}
{
    \zA{在袖里。}
}
{
    \eA{[English source not present in the checked-in Bencraft/Joly files.]}
}
{
}

\Verse{116}
{
    \zA{心内十分 罕异，揣摩此物从何而来，且不形于声色，到了迎春房里。}
}
{
    \eA{[English source not present in the checked-in Bencraft/Joly files.]}
}
{
}

\Verse{117}
{
    \zA{迎春正因他乳母获罪， 心中不自在，忽报母亲来了，遂接入。}
}
{
    \eA{[English source not present in the checked-in Bencraft/Joly files.]}
}
{
}

\Verse{118}
{
    \zA{奉茶毕，邢夫人因说道：“你这么大了，你 那奶妈子行此事，你也不说说他。}
}
{
    \eA{[English source not present in the checked-in Bencraft/Joly files.]}
}
{
}

\Verse{119}
{
    \zA{如今别人都好好的，偏咱们的人做出这事来，什 么意思？”}
}
{
    \eA{[English source not present in the checked-in Bencraft/Joly files.]}
}
{
}

\Verse{120}
{
    \zA{迎春低头弄衣带，半晌答道：“我说他两次，他不听，也叫我没法儿。}
}
{
    \eA{[English source not present in the checked-in Bencraft/Joly files.]}
}
{
}

\Verse{121}
{
    \zA{况因他是妈妈，只有他说我的，没有我说他的。”}
}
{
    \eA{[English source not present in the checked-in Bencraft/Joly files.]}
}
{
}

\Verse{122}
{
    \zA{邢夫人道：“胡说。}
}
{
    \eA{[English source not present in the checked-in Bencraft/Joly files.]}
}
{
}

\Verse{123}
{
    \zA{你不好了， 他原该说；如今他犯了法，你就该拿出姑娘的身分来。}
}
{
    \eA{[English source not present in the checked-in Bencraft/Joly files.]}
}
{
}

\Verse{124}
{
    \zA{他敢不依，你就回我去才是。}
}
{
    \eA{[English source not present in the checked-in Bencraft/Joly files.]}
}
{
}

\Verse{125}
{
    \zA{如今直等外人共知，这可是什么意思!再者：放头儿，还只怕他巧语花言的和你借 贷些簪环衣裳做本钱。}
}
{
    \eA{[English source not present in the checked-in Bencraft/Joly files.]}
}
{
}

\Verse{126}
{
    \zA{你这心活面软，未必不周济他些。}
}
{
    \eA{[English source not present in the checked-in Bencraft/Joly files.]}
}
{
}

\Verse{127}
{
    \zA{若被他骗了去，我是一个 钱没有的，看你明日怎么过节？”}
}
{
    \eA{[English source not present in the checked-in Bencraft/Joly files.]}
}
{
}

\Verse{128}
{
    \zA{迎春不语，只低着头。}
}
{
    \eA{[English source not present in the checked-in Bencraft/Joly files.]}
}
{
}

\Verse{129}
{
    \zA{邢夫人见他这般，因冷笑 道：“你是大老爷跟前的人养的，这里探丫头是二老爷跟前的人养的，出身一样， 你娘比赵姨娘强十分，你也该比探丫头强才是。}
}
{
    \eA{[English source not present in the checked-in Bencraft/Joly files.]}
}
{
}

\Verse{130}
{
    \zA{怎么你反不及他一点？}
}
{
    \eA{[English source not present in the checked-in Bencraft/Joly files.]}
}
{
}

\Verse{131}
{
    \zA{倒是我无儿 女的一生干净，也不能惹人笑话！”}
}
{
    \eA{[English source not present in the checked-in Bencraft/Joly files.]}
}
{
}

\Verse{132}
{
    \zA{人回：“琏二奶奶来了。”}
}
{
    \eA{[English source not present in the checked-in Bencraft/Joly files.]}
}
{
}

\Verse{133}
{
    \zA{邢夫人听了，冷笑 两声，命人出去说：“请他自己养病，我这里不用他伺候。”}
}
{
    \eA{[English source not present in the checked-in Bencraft/Joly files.]}
}
{
}

\Verse{134}
{
    \zA{接着又有探事的小丫 头来报说：“老太太醒了。”}
}
{
    \eA{[English source not present in the checked-in Bencraft/Joly files.]}
}
{
}

\Verse{135}
{
    \zA{邢夫人方起身往前边来。}
}
{
    \eA{[English source not present in the checked-in Bencraft/Joly files.]}
}
{
}

\Verse{136}
{
    \zA{迎春送至院外方回。}
}
{
    \eA{[English source not present in the checked-in Bencraft/Joly files.]}
}
{
}

\Verse{137}
{
    \zA{绣橘因说道：“如何？}
}
{
    \eA{[English source not present in the checked-in Bencraft/Joly files.]}
}
{
}

\Verse{138}
{
    \zA{前儿我回姑娘：‘那一个攒珠累金 凤，竟不知那里去了。’}
}
{
    \eA{[English source not present in the checked-in Bencraft/Joly files.]}
}
{
}

\Verse{139}
{
    \zA{回了姑娘，竟不问一声儿。}
}
{
    \eA{[English source not present in the checked-in Bencraft/Joly files.]}
}
{
}

\Verse{140}
{
    \zA{我说：‘必是老奶奶拿去当了 银子放头儿了。’}
}
{
    \eA{[English source not present in the checked-in Bencraft/Joly files.]}
}
{
}

\Verse{141}
{
    \zA{姑娘不信，只说司棋收着，叫问司棋。}
}
{
    \eA{[English source not present in the checked-in Bencraft/Joly files.]}
}
{
}

\Verse{142}
{
    \zA{司棋虽病，心里却明白， 说：‘没有收起来，还在书架上匣里放着，预备八月十五要戴呢。’}
}
{
    \eA{[English source not present in the checked-in Bencraft/Joly files.]}
}
{
}

\Verse{143}
{
    \zA{姑娘该叫人去 问老奶奶一声。”}
}
{
    \eA{[English source not present in the checked-in Bencraft/Joly files.]}
}
{
}

\Verse{144}
{
    \zA{迎春道：“何用问？}
}
{
    \eA{[English source not present in the checked-in Bencraft/Joly files.]}
}
{
}

\Verse{145}
{
    \zA{那自然是他拿了去摘了肩儿了。}
}
{
    \eA{[English source not present in the checked-in Bencraft/Joly files.]}
}
{
}

\Verse{146}
{
    \zA{我只说他悄 悄的拿了出去，不过一时半晌，仍旧悄悄的放在里头，谁知他就忘了。}
}
{
    \eA{[English source not present in the checked-in Bencraft/Joly files.]}
}
{
}

\Verse{147}
{
    \zA{今日偏又闹 出来，问他也无益。”}
}
{
    \eA{[English source not present in the checked-in Bencraft/Joly files.]}
}
{
}

\Verse{148}
{
    \zA{绣橘道：“何曾是忘记？}
}
{
    \eA{[English source not present in the checked-in Bencraft/Joly files.]}
}
{
}

\Verse{149}
{
    \zA{他是试准了姑娘的性格儿才这么着。}
}
{
    \eA{[English source not present in the checked-in Bencraft/Joly files.]}
}
{
}

\Verse{150}
{
    \zA{如今我有个主意：到二奶奶屋里，将此事回了，他或着人要，他或省事拿几吊钱来 替他赎了，如何？”}
}
{
    \eA{[English source not present in the checked-in Bencraft/Joly files.]}
}
{
}

\Verse{151}
{
    \zA{迎春忙道：“罢，罢，省事些好。}
}
{
    \eA{[English source not present in the checked-in Bencraft/Joly files.]}
}
{
}

\Verse{152}
{
    \zA{宁可没有了，又何必生事？”}
}
{
    \eA{[English source not present in the checked-in Bencraft/Joly files.]}
}
{
}

\Verse{153}
{
    \zA{绣橘道：“姑娘怎么这样软弱？}
}
{
    \eA{[English source not present in the checked-in Bencraft/Joly files.]}
}
{
}

\Verse{154}
{
    \zA{都要省起事来，将来连姑娘还骗了去。}
}
{
    \eA{[English source not present in the checked-in Bencraft/Joly files.]}
}
{
}

\Verse{155}
{
    \zA{我竟去的是。”}
}
{
    \eA{[English source not present in the checked-in Bencraft/Joly files.]}
}
{
}

\Verse{156}
{
    \zA{说着便走。}
}
{
    \eA{[English source not present in the checked-in Bencraft/Joly files.]}
}
{
}

\Verse{157}
{
    \zA{迎春便不言语，只好由他。}
}
{
    \eA{[English source not present in the checked-in Bencraft/Joly files.]}
}
{
}

\Verse{158}
{
    \zA{谁知迎春的乳母之媳玉柱儿媳妇为他婆婆得罪，来求迎春去讨情，他们正说金 凤一事，且不进去。}
}
{
    \eA{[English source not present in the checked-in Bencraft/Joly files.]}
}
{
}

\Verse{159}
{
    \zA{也因素日迎春懦弱，他们都不放在心上；如今见绣橘立意去回 凤姐，又看这事脱不过去，只得进来，陪笑先向绣橘说：“姑娘，你别去生事。}
}
{
    \eA{[English source not present in the checked-in Bencraft/Joly files.]}
}
{
}

\Verse{160}
{
    \zA{姑 娘的金丝凤，原是我们老奶奶老糊涂了，输了几个钱，没的捞梢，所以借去，不想 今日弄出事来。}
}
{
    \eA{[English source not present in the checked-in Bencraft/Joly files.]}
}
{
}

\Verse{161}
{
    \zA{虽然这样，到底主子的东西，我们不敢迟误，终久是要赎的。}
}
{
    \eA{[English source not present in the checked-in Bencraft/Joly files.]}
}
{
}

\Verse{162}
{
    \zA{如今 还要求姑娘看着从小儿吃奶的情，往老太太那边去讨一个情儿，救出他来才好。”}
}
{
    \eA{[English source not present in the checked-in Bencraft/Joly files.]}
}
{
}

\Verse{163}
{
    \zA{迎春便说道：“好嫂子，你趁早打了这妄想。}
}
{
    \eA{[English source not present in the checked-in Bencraft/Joly files.]}
}
{
}

\Verse{164}
{
    \zA{要等我去说情儿，等到明年，也是不 中用的。}
}
{
    \eA{[English source not present in the checked-in Bencraft/Joly files.]}
}
{
}

\Verse{165}
{
    \zA{方才连宝姐姐林妹妹，大伙儿说情，老太太还不依，何况是我一个人？}
}
{
    \eA{[English source not present in the checked-in Bencraft/Joly files.]}
}
{
}

\Verse{166}
{
    \zA{我 自己臊还臊不过来，还去讨臊去？”}
}
{
    \eA{[English source not present in the checked-in Bencraft/Joly files.]}
}
{
}

\Verse{167}
{
    \zA{绣橘便说：“赎金凤是一件事，说情是一件事， 别绞在一处。}
}
{
    \eA{[English source not present in the checked-in Bencraft/Joly files.]}
}
{
}

\Verse{168}
{
    \zA{难道姑娘不去说情，你就不赔了不成？}
}
{
    \eA{[English source not present in the checked-in Bencraft/Joly files.]}
}
{
}

\Verse{169}
{
    \zA{嫂子且取了金凤来再说。”}
}
{
    \eA{[English source not present in the checked-in Bencraft/Joly files.]}
}
{
}

\Verse{170}
{
    \zA{玉 柱儿家的听见迎春如此拒绝他，绣橘的话又锋利，无可回答，一时脸上过不去，也 明欺迎春素日好性儿，乃向绣橘说道：“姑娘，你别太张势了!你满家子算一算，
        谁的妈妈奶奶不仗着主子哥儿姐儿得些便宜，偏咱们就这样‘丁是丁，卯是卯’的？}
}
{
    \eA{[English source not present in the checked-in Bencraft/Joly files.]}
}
{
}

\Verse{171}
{
    \zA{只许你们偷偷摸摸的哄骗了去。}
}
{
    \eA{[English source not present in the checked-in Bencraft/Joly files.]}
}
{
}
\ChapterImage{img/chapters/142第73-74回 因春囊重托王善保.jpg}


\Verse{172}
{
    \zA{自从邢姑娘来了，太太吩咐一个月俭省出一两银子 来给舅太太去，这里饶添了邢姑娘的使费，反少了一两银子。}
}
{
    \eA{[English source not present in the checked-in Bencraft/Joly files.]}
}
{
}

\Verse{173}
{
    \zA{时常短了这个，少了 那个，那不是我们供给？}
}
{
    \eA{[English source not present in the checked-in Bencraft/Joly files.]}
}
{
}

\Verse{174}
{
    \zA{谁又要去？}
}
{
    \eA{[English source not present in the checked-in Bencraft/Joly files.]}
}
{
}

\Verse{175}
{
    \zA{不过大家将就些罢了。}
}
{
    \eA{[English source not present in the checked-in Bencraft/Joly files.]}
}
{
}

\Verse{176}
{
    \zA{算到今日少说也有三十两 了，我们这一向的钱岂不白填了限呢？”}
}
{
    \eA{[English source not present in the checked-in Bencraft/Joly files.]}
}
{
}

\Verse{177}
{
    \zA{绣橘不待说完，便啐了一口，道：“做什 么你白填了三十两？}
}
{
    \eA{[English source not present in the checked-in Bencraft/Joly files.]}
}
{
}

\Verse{178}
{
    \zA{我且和你算算账!姑娘要了些什么东西？”}
}
{
    \eA{[English source not present in the checked-in Bencraft/Joly files.]}
}
{
}

\Verse{179}
{
    \zA{迎春听了这媳妇发邢 夫人之私意，忙止道：“罢，罢!不能拿了金凤来，你不必拉三扯四的乱嚷。}
}
{
    \eA{[English source not present in the checked-in Bencraft/Joly files.]}
}
{
}

\Verse{180}
{
    \zA{我也 不要那凤了。}
}
{
    \eA{[English source not present in the checked-in Bencraft/Joly files.]}
}
{
}

\Verse{181}
{
    \zA{就是太太问时，我只说丢了，也妨碍不着你什么，你出去歇歇儿去罢。}
}
{
    \eA{[English source not present in the checked-in Bencraft/Joly files.]}
}
{
}

\Verse{182}
{
    \zA{何苦呢？”}
}
{
    \eA{[English source not present in the checked-in Bencraft/Joly files.]}
}
{
}

\Verse{183}
{
    \zA{一面叫绣橘倒茶来。}
}
{
    \eA{[English source not present in the checked-in Bencraft/Joly files.]}
}
{
}

\Verse{184}
{
    \zA{绣橘又气又急，因说道：“姑娘虽不怕，我是做什 么的？}
}
{
    \eA{[English source not present in the checked-in Bencraft/Joly files.]}
}
{
}

\Verse{185}
{
    \zA{把姑娘的东西丢了，他倒赖说姑娘使了他的钱，这如今竟要准折起来。}
}
{
    \eA{[English source not present in the checked-in Bencraft/Joly files.]}
}
{
}

\Verse{186}
{
    \zA{倘或 太太问姑娘为什么使了这些钱，敢是我们就中取势？}
}
{
    \eA{[English source not present in the checked-in Bencraft/Joly files.]}
}
{
}

\Verse{187}
{
    \zA{这还了得！”}
}
{
    \eA{[English source not present in the checked-in Bencraft/Joly files.]}
}
{
}

\Verse{188}
{
    \zA{一行说，一行就 哭了。}
}
{
    \eA{[English source not present in the checked-in Bencraft/Joly files.]}
}
{
}

\Verse{189}
{
    \zA{司棋听不过，只得勉强过来，帮着绣橘问着那媳妇。}
}
{
    \eA{[English source not present in the checked-in Bencraft/Joly files.]}
}
{
}

\Verse{190}
{
    \zA{迎春劝止不住，自拿了 一本《太上感应篇》去看。}
}
{
    \eA{[English source not present in the checked-in Bencraft/Joly files.]}
}
{
}

\Verse{191}
{
    \zA{三人正没开交，可巧宝钗、黛玉、宝琴、探春等，因恐迎春今日不自在，都约 着来安慰。}
}
{
    \eA{[English source not present in the checked-in Bencraft/Joly files.]}
}
{
}

\Verse{192}
{
    \zA{他们走至院中，听见几个人讲究，探春从纱窗内一看，只见迎春倚在床 上看书，若有不闻之状，探春也笑了。}
}
{
    \eA{[English source not present in the checked-in Bencraft/Joly files.]}
}
{
}

\Verse{193}
{
    \zA{小丫头们忙打起帘子报道：“姑娘们来了。”}
}
{
    \eA{[English source not present in the checked-in Bencraft/Joly files.]}
}
{
}

\Verse{194}
{
    \zA{迎春放下书起身。}
}
{
    \eA{[English source not present in the checked-in Bencraft/Joly files.]}
}
{
}

\Verse{195}
{
    \zA{那媳妇见有人来，且又有探春在内，不劝自止了，遂趁便就走。}
}
{
    \eA{[English source not present in the checked-in Bencraft/Joly files.]}
}
{
}

\Verse{196}
{
    \zA{探春坐下，便问：“才刚谁在这里说话，倒像拌嘴似的？”}
}
{
    \eA{[English source not present in the checked-in Bencraft/Joly files.]}
}
{
}

\Verse{197}
{
    \zA{迎春笑道：“没有什么， 左不过他们小题大做罢了，何必问他？”}
}
{
    \eA{[English source not present in the checked-in Bencraft/Joly files.]}
}
{
}

\Verse{198}
{
    \zA{探春笑道：“我才听见什么‘金凤’，又 是什么‘没有钱，只合我们奴才要’。}
}
{
    \eA{[English source not present in the checked-in Bencraft/Joly files.]}
}
{
}

\Verse{199}
{
    \zA{谁和奴才要钱了？}
}
{
    \eA{[English source not present in the checked-in Bencraft/Joly files.]}
}
{
}

\Verse{200}
{
    \zA{难道姐姐和奴才要钱不成？”}
}
{
    \eA{[English source not present in the checked-in Bencraft/Joly files.]}
}
{
}

\Verse{201}
{
    \zA{司棋绣橘道：“姑娘说的是了!姑娘何曾和他要什么了？”}
}
{
    \eA{[English source not present in the checked-in Bencraft/Joly files.]}
}
{
}

\Verse{202}
{
    \zA{探春笑道：“姐姐既没 有和他要，必定是我们和他们要了不成？}
}
{
    \eA{[English source not present in the checked-in Bencraft/Joly files.]}
}
{
}

\Verse{203}
{
    \zA{你叫他进来，我倒要问问他。”}
}
{
    \eA{[English source not present in the checked-in Bencraft/Joly files.]}
}
{
}

\Verse{204}
{
    \zA{迎春笑道： “这话又可笑。}
}
{
    \eA{[English source not present in the checked-in Bencraft/Joly files.]}
}
{
}

\Verse{205}
{
    \zA{你们又无沾碍，何必如此？”}
}
{
    \eA{[English source not present in the checked-in Bencraft/Joly files.]}
}
{
}

\Verse{206}
{
    \zA{探春道：“这倒不然。}
}
{
    \eA{[English source not present in the checked-in Bencraft/Joly files.]}
}
{
}

\Verse{207}
{
    \zA{我和姐姐一样。}
}
{
    \eA{[English source not present in the checked-in Bencraft/Joly files.]}
}
{
}

\Verse{208}
{
    \zA{姐姐的事，和我一般。}
}
{
    \eA{[English source not present in the checked-in Bencraft/Joly files.]}
}
{
}

\Verse{209}
{
    \zA{他说姐姐，即是说我；我那边有人怨我，姐姐听见，也是合 怨姐姐一样。}
}
{
    \eA{[English source not present in the checked-in Bencraft/Joly files.]}
}
{
}

\Verse{210}
{
    \zA{咱们是主子，自然不理论那些钱财小事，只知想起什么要什么，也是 有的事。}
}
{
    \eA{[English source not present in the checked-in Bencraft/Joly files.]}
}
{
}

\Verse{211}
{
    \zA{但不知累丝凤怎么又夹在里头？”}
}
{
    \eA{[English source not present in the checked-in Bencraft/Joly files.]}
}
{
}

\Verse{212}
{
    \zA{那玉柱儿媳妇生恐绣橘等告出他来，遂 忙进来用话掩饰。}
}
{
    \eA{[English source not present in the checked-in Bencraft/Joly files.]}
}
{
}

\Verse{213}
{
    \zA{探春深知其意，因笑道：“你们所以糊涂!如今你奶奶已得了不 是，趁此求二奶奶，把方才的钱未曾散人的拿出些来赎来就完了。}
}
{
    \eA{[English source not present in the checked-in Bencraft/Joly files.]}
}
{
}

\Verse{214}
{
    \zA{比不得没闹出来， 大家都藏着留脸面。}
}
{
    \eA{[English source not present in the checked-in Bencraft/Joly files.]}
}
{
}

\Verse{215}
{
    \zA{如今既是没了脸，趁此时，总有十个罪也只一人受罚，没有砍 两颗头的理。}
}
{
    \eA{[English source not present in the checked-in Bencraft/Joly files.]}
}
{
}

\Verse{216}
{
    \zA{你依我说，竟是和二奶奶趁便说去。}
}
{
    \eA{[English source not present in the checked-in Bencraft/Joly files.]}
}
{
}

\Verse{217}
{
    \zA{在这里大声小气，如何使得！”}
}
{
    \eA{[English source not present in the checked-in Bencraft/Joly files.]}
}
{
}

\Verse{218}
{
    \zA{这媳妇被探春说出真病，也无可赖了，只不敢往凤姐处自首。}
}
{
    \eA{[English source not present in the checked-in Bencraft/Joly files.]}
}
{
}

\Verse{219}
{
    \zA{探春笑道：“我不听 见便罢，既听见，少不得替你们分解分解。”}
}
{
    \eA{[English source not present in the checked-in Bencraft/Joly files.]}
}
{
}

\Verse{220}
{
    \zA{谁知探春早使了眼色与侍书，侍书出去了。}
}
{
    \eA{[English source not present in the checked-in Bencraft/Joly files.]}
}
{
}

\Verse{221}
{
    \zA{这里正说话，忽见平儿进来。}
}
{
    \eA{[English source not present in the checked-in Bencraft/Joly files.]}
}
{
}

\Verse{222}
{
    \zA{宝琴 拍手笑道：“三姐姐敢是有驱神召将的符术？”}
}
{
    \eA{[English source not present in the checked-in Bencraft/Joly files.]}
}
{
}

\Verse{223}
{
    \zA{黛玉笑道：“这倒不是道家法术， 倒是用兵最精的所谓‘守如处女，出如脱兔’，‘出其不备’的妙策。”}
}
{
    \eA{[English source not present in the checked-in Bencraft/Joly files.]}
}
{
}

\Verse{224}
{
    \zA{二人取笑， 宝钗便使眼色与二人，遂以别话岔开。}
}
{
    \eA{[English source not present in the checked-in Bencraft/Joly files.]}
}
{
}

\Verse{225}
{
    \zA{探春见平儿来了，遂问：“你奶奶可好些了？}
}
{
    \eA{[English source not present in the checked-in Bencraft/Joly files.]}
}
{
}

\Verse{226}
{
    \zA{真是病糊涂了，事事都不在心上，叫我们受这样委屈。”}
}
{
    \eA{[English source not present in the checked-in Bencraft/Joly files.]}
}
{
}

\Verse{227}
{
    \zA{平儿忙道：“谁敢给姑娘 气受？}
}
{
    \eA{[English source not present in the checked-in Bencraft/Joly files.]}
}
{
}

\Verse{228}
{
    \zA{姑娘吩咐我。”}
}
{
    \eA{[English source not present in the checked-in Bencraft/Joly files.]}
}
{
}

\Verse{229}
{
    \zA{那玉柱儿媳妇方慌了手脚，遂上来赶着平儿叫：“姑娘坐下， 让我说原故，姑娘请听。”}
}
{
    \eA{[English source not present in the checked-in Bencraft/Joly files.]}
}
{
}

\Verse{230}
{
    \zA{平儿正色道：“姑娘这里说话，也有你混插嘴的理吗! 你但凡知礼，该在外头伺候，也有外头的媳妇们无故到姑娘屋里来的？”}
}
{
    \eA{[English source not present in the checked-in Bencraft/Joly files.]}
}
{
}

\Verse{231}
{
    \zA{绣橘道： “你不知我们这屋里是没礼的，谁爱来就来。”}
}
{
    \eA{[English source not present in the checked-in Bencraft/Joly files.]}
}
{
}

\Verse{232}
{
    \zA{平儿道：“都是你们不是!姑娘好 性儿，你们就该打出去，然后再回太太去才是。”}
}
{
    \eA{[English source not present in the checked-in Bencraft/Joly files.]}
}
{
}

\Verse{233}
{
    \zA{柱儿媳妇见平儿出了言，红了脸， 才退出去。}
}
{
    \eA{[English source not present in the checked-in Bencraft/Joly files.]}
}
{
}

\Verse{234}
{
    \zA{探春接着道：“我且告诉你：要是别人得罪了我，倒还罢了。}
}
{
    \eA{[English source not present in the checked-in Bencraft/Joly files.]}
}
{
}

\Verse{235}
{
    \zA{如今这柱 儿媳妇和他婆婆，仗着是嬷嬷，又瞅着二姐姐好性儿，私自拿了首饰去赌钱，而且
        还捏造假账，逼着去讨情，和这两个丫头在卧房里大嚷大叫，二姐姐竟不能辖治。}
}
{
    \eA{[English source not present in the checked-in Bencraft/Joly files.]}
}
{
}

\Verse{236}
{
    \zA{所以我看不过，才请你来问一声：还是他本是天外的人，不知道理？}
}
{
    \eA{[English source not present in the checked-in Bencraft/Joly files.]}
}
{
}

\Verse{237}
{
    \zA{还是有谁主使 他如此，先把二姐姐制伏了，然后就要治我和四姑娘了？”}
}
{
    \eA{[English source not present in the checked-in Bencraft/Joly files.]}
}
{
}

\Verse{238}
{
    \zA{平儿忙陪笑道：“姑娘 怎么今日说出这话来？}
}
{
    \eA{[English source not present in the checked-in Bencraft/Joly files.]}
}
{
}

\Verse{239}
{
    \zA{我们奶奶如何担得起！”}
}
{
    \eA{[English source not present in the checked-in Bencraft/Joly files.]}
}
{
}

\Verse{240}
{
    \zA{探春冷笑道：“俗语说的，‘物伤 其类，唇亡齿寒’，我自然有些心惊么。”}
}
{
    \eA{[English source not present in the checked-in Bencraft/Joly files.]}
}
{
}

\Verse{241}
{
    \zA{平儿问迎春道：“若论此事，本好处的。}
}
{
    \eA{[English source not present in the checked-in Bencraft/Joly files.]}
}
{
}

\Verse{242}
{
    \zA{但只他是姑娘的奶嫂，姑娘怎么样呢？”}
}
{
    \eA{[English source not present in the checked-in Bencraft/Joly files.]}
}
{
}

\Verse{243}
{
    \zA{当下迎春只合宝钗看《感应篇》故事，究竟连探春的话也没听见，忽见平儿如此说， 仍笑道：“问我，我也没什么法子。}
}
{
    \eA{[English source not present in the checked-in Bencraft/Joly files.]}
}
{
}

\Verse{244}
{
    \zA{他们的不是，自作自受，我也不能讨情，我也 不去加责，就是了。}
}
{
    \eA{[English source not present in the checked-in Bencraft/Joly files.]}
}
{
}

\Verse{245}
{
    \zA{至于私自拿去的东西，送来我收下，不送来我也不要了。}
}
{
    \eA{[English source not present in the checked-in Bencraft/Joly files.]}
}
{
}

\Verse{246}
{
    \zA{太太 们要来问我，可以隐瞒遮饰的过去，是他的造化；要瞒不住我也没法儿，没有个为 他们反欺枉太太们的理，少不得直说。}
}
{
    \eA{[English source not present in the checked-in Bencraft/Joly files.]}
}
{
}

\Verse{247}
{
    \zA{你们要说我好性儿，没个决断；有好主意可 以八面周全，不叫太太们生气，任凭你们处治，我也不管。”}
}
{
    \eA{[English source not present in the checked-in Bencraft/Joly files.]}
}
{
}

\Verse{248}
{
    \zA{众人听了，都好笑起 来。}
}
{
    \eA{[English source not present in the checked-in Bencraft/Joly files.]}
}
{
}

\Verse{249}
{
    \zA{黛玉笑道：“真是‘虎狼屯于阶陛，尚谈因果’。}
}
{
    \eA{[English source not present in the checked-in Bencraft/Joly files.]}
}
{
}

\Verse{250}
{
    \zA{要是二姐姐是个男人，一家 上下这些人，又如何裁治他们？”}
}
{
    \eA{[English source not present in the checked-in Bencraft/Joly files.]}
}
{
}

\Verse{251}
{
    \zA{迎春笑道：“正是，多少男人衣租食税，及至事 到临头，尚且如此。}
}
{
    \eA{[English source not present in the checked-in Bencraft/Joly files.]}
}
{
}

\Verse{252}
{
    \zA{况且‘太上’说的好，救人急难，最是阴骘事。}
}
{
    \eA{[English source not present in the checked-in Bencraft/Joly files.]}
}
{
}

\Verse{253}
{
    \zA{我虽不能救人， 何苦来白白去和人结怨结仇，作那样无益有损的事呢？”}
}
{
    \eA{[English source not present in the checked-in Bencraft/Joly files.]}
}
{
}

\Verse{254}
{
    \zA{一语未了，只听又有一人 来了。}
}
{
    \eA{[English source not present in the checked-in Bencraft/Joly files.]}
}
{
}

\Verse{255}
{
    \zA{不知是谁，下回分解。}
}
{
    \eA{[English source not present in the checked-in Bencraft/Joly files.]}
}
{
}

\Verse{256}
{
    \zA{(本章完)}
}
{
    \eA{[English source not present in the checked-in Bencraft/Joly files.]}
}
{
}
